\subsection{Differentiable functions at a point}

1.2.1. Let \(f\) be a function defined in a neighborhood of the point \(x_{0}\) of E and with values in F. We say that \(f\) is differentiable at \(x_{0}\) if there exists a continuous affine function \(v\) from E into F having at \(x_{0}\) a contact of order \(\geqslant 1\) with \(f\) . This mapping \(v\) is unique; there exists one and only one continuous linear mapping, denoted \(Df(x_{0})\) , from E into F such that:

\[
v (x) = v \left(x _ {0}\right) + D f \left(x _ {0}\right). \left(x - x _ {0}\right).
\]

If one chooses a norm on E, this is equivalent to:

\(f(x_0 + h) \equiv f(x_0) + Df(x_0).h \mod o(\|h\|)\) for \(h\) tending toward 0,

which can also be written in the form:

\[
\lim _ {h \rightarrow 0, h \neq 0} \frac {\| f (x _ {0} + h) - f (x _ {0}) - D f (x _ {0}) . h \| _ {\gamma}}{\| h \|} = 0
\]

for every continuous seminorm γ on F.

The element \(Df(x_{0})\) of \(\mathcal{L}(\mathrm{E},\mathrm{F})\) is called the derivative of f at \(x_{0}\) . One sometimes writes \(D_{h}f(x_{0})\) for \(Df(x_{0}).h\) . It is an element of F defined by the relation:

\[
\mathrm{D} _ {h} f (x _ {0}) = \lim _ {t \rightarrow 0, t \neq 0} \frac {f (x _ {0} + t h) - f (x _ {0})}{t}.
\]

1.2.2. We say that a function \(f\) is strictly differentiable at \(x_{0}\) if it is differentiable at \(x_{0}\) and if, for every norm defining the topology of E, one has the relation:

\[
f (y) - f (z) \equiv D f (x _ {0}). (y - z) \quad \text {   mod   } o (\| y - z \|)
\]

for \((y, z)\) tending toward \((x_0, x_0)\) in \(E \times E\) . For this, it suffices that this condition be satisfied for one norm defining the topology of \(E\) . Suppose moreover \(E\) and \(F\) normed; for every number \(c > \|Df(x_0)\|\) there then exists a neighborhood V of \(x_{0}\) such that \(\|f(y)-f(z)\|\leqslant c.\|y-z\|\) for y, z in V; this implies that f is uniformly continuous in V.

1.2.3. The property of a function \(f\) being differentiable or strictly differentiable at \(x_{0}\) depends only on the germ of \(f\) at \(x_{0}\) . The germs of differentiable functions at \(x_{0}\) form a vector subspace \(\mathcal{V}\) of the space of all germs and the mapping \(f\mapsto\mathrm{D}f(x_{0})\) of \(\mathcal{V}\) to \(\mathcal{L}(\mathrm{E};\mathrm{F})\) is linear. The germs of functions strictly differentiable at \(x_{0}\) form a vector subspace of \(\mathcal{V}\) .

1.2.4. A function differentiable at \(x_0\) is continuous at \(x_0\) .

1.2.5. When E = K, the map \(u \mapsto u(1)\) is an isomorphism of \(\mathcal{L}(\mathbf{E}; \mathbf{F})\) onto F; if the function f is differentiable at \(x_{0}\) , the element

\[
f ^ {\prime} (x _ {0}) = \mathrm{D} f (x _ {0}). 1
\]

is none other than the derivative of \(f\) at \(x_0\) in the sense given in Fonct. Var. Réelle, chap. I, § 1, n° 6, remark 2.

